%\needspace{4\baselineskip}
\color{unidarkgreen}{\subsection{Направленная высокочастотная защита}}
\color{black}

\begin{enumerate}[label=\arabic{section}.\arabic{subsection}.\arabic{enumi}, labelsep=4pt, leftmargin=0pt, itemindent=57pt, itemsep=0pt, parsep=5pt] % если что можно крутить itemsep

%%%%%%%%%%%%%%%%%%%%%%%%%%%%%%%%%%%%%%%%%%%%%%%%%% ОБШИЕ СВЕДЕНИЯ %%%%%%%%%%%%%%%%%%%%%%%%%%%%%%%%%%%%%%%%%%%%%%%%%%
\item Общие сведения

\begin{enumerate}[label=\arabic{section}.\arabic{subsection}.\arabic{enumi}.\arabic*, labelsep=4pt, leftmargin=0em, itemindent=65pt, parsep=0pt]

\item
Направленная высокочастотная защита (НВЧЗ) является быстродействующей защитой линии с абсолютной селективностью и выполняет функцию основной защиты линии от всех видов КЗ. 

\item
НВЧЗ реализуется в виде полукомплектов, устанавливаемых на концах линии, каждый из которых включает микропроцессорный терминал релейной защиты и комплект оборудования высокочастотной (ВЧ) связи, обеспечивающий формирование и приём ВЧ-сигналов по каналу связи с использованием внешних приёмопередатчиков (ПП).

\item 
НВЧЗ селективно срабатывает:
\begin{itemize}
    \item при всех видах КЗ, в том числе при постановке линии под напряжение; 
    \item при переходе внешнего КЗ во внутреннее, в том числе при реверсе мощности.
\end{itemize}

\item 
Обеспечивается несрабатывание защиты:
\begin{itemize} 
\item при постановке линии под напряжение и включении в транзит при отсутствии КЗ;
\item при внешних КЗ, в том числе при реверсе мощности;
\item при синхронных качаниях и в режиме асинхронного хода;
\item при КЗ за трансформаторами ответвлений (отпаек) на линиях с ответвлениями;
\item при бросках тока намагничивания трансформаторов без КЗ;
\item при несимметрии токов, вызванной тяговой нагрузкой;
\item при внешних КЗ с насыщением трансформаторов тока (при соблюдении заявленных производителем требований к ТТ).
\end{itemize}

\item
Принцип действия НВЧЗ основан на сравнении направлений мощности по концам защищаемой линии, определяемых органом направления мощности обратной последовательности (ОНМОП) при несимметричных КЗ и направленным реле сопротивления при симметричных.
Сравнение направлений мощности осуществляется косвенно – посредством обмена ВЧ-сигналами между полукомплектами. 

\item
НВЧЗ содержит две основные группы измерительных органов (ИО): блокирующие – действующие на пуск ВЧ-передатчика, и отключающие – действующие на отключение выключателя.
В начальный момент повреждения срабатывают блокирующие органы, обеспечивая ускоренный запуск ВЧ-передатчика и передачу блокирующего ВЧ-сигнала на противоположный конец линии.
Место КЗ, внутри или вне защищаемой линии, определяется по отсутствию или наличию постоянного ВЧ-сигнала в канале связи. 

\item
Если повреждение находится в пределах защищаемой линии, то после срабатывания отключающих органов, имеющих более грубые уставки, ВЧ-передатчики каждого полукомплекта останавливаются, при этом блокирующий ВЧ-сигнал снимается, и оба полукомплекта действуют на отключение своих выключателей.
Если повреждение находится вне защищаемой линии, «за спиной» одного из полукомплектов, то в этом полукомплекте отключающие органы не срабатывают, его ВЧ-передатчик остается в работе, и блокирующий ВЧ-сигнал продолжает передаваться по каналу связи, предотвращая срабатывание защиты обоих полукомплектов.

\end{enumerate}


%%%%%%%%%%%%%%%%%%%%%%%%%%%%%%%%%%%%%%%%%%%%%%%%%% БЛОКИРУЮЩИЕ ИО %%%%%%%%%%%%%%%%%%%%%%%%%%%%%%%%%%%%%%%%%%%%%%%%%%
\item
Блокирующие измерительные органы

\begin{enumerate}[label=\arabic{section}.\arabic{subsection}.\arabic{enumi}.\arabic*, labelsep=4pt, leftmargin=0em, itemindent=65pt, parsep=0pt]

\item
Команда на пуск ВЧ-передатчика (логический сигнал \textit{«Пуск ВЧ от НВЧЗ»}) формируется при срабатывании блокирующих органов, отстроенных от нагрузочного режима.
Блокирующие органы реализованы с возможностью независимой настройки их уставок по:
\begin{itemize} 
\item току обратной последовательности \textit{«ИО I2 блок НВЧЗ»};
\item приращению тока прямой последовательности \textit{«ИО dI1 блок НВЧЗ»};
\item приращению тока обратной последовательности \textit{«ИО dI2 блок НВЧЗ»};
\item напряжению обратной последовательности \textit{«ИО U2 блок НВЧЗ»};
\item сопротивлению «РС блок НВЧЗ».
\end{itemize}

\item
Блокирующие органы по приращениям токов прямой и обратной последовательностей рекомендуется применять в сетях с тяговой нагрузкой, для которых характерны значительные небалансы токов обратной последовательности.
Время возврата этих органов составляет 250 мс. 

\item
Блокирующий орган по сопротивлению (блокирующее РС) имеет обратнонаправленную характеристику срабатывания с охватом начала координат, которая задается уставками \textbf{«Rср РС блок НВЧЗ»}, \textbf{«Хср РС блок НВЧЗ»} и \textbf{«Ф1 РС блок НВЧЗ»}.
Для отстройки от нагрузки предусмотрен вырез в характеристике срабатывания, который определяется уставками \textbf{«Rнг РС блок НВЧЗ»} и \textbf{«Фнг РС блок НВЧЗ»}.
Работа блокирующего РС может быть заблокирована при выявлении неисправности в цепях напряжения при введенном программном ключе \textbf{«БНН РС блок НВЧЗ»}.

\medskip

\begin{figure}[H]
\centering
\includegraphics[width=0.5\textwidth,height=0.5\textheight,keepaspectratio]{\fbpath/561.0771 НВЧЗ/img/RS_NVChZ.pdf}
\captionof{figure}{Характеристика срабатывания блокирующего РС НВЧЗ}\label{pict_nvchz:RS_NVChZ}
\end{figure}

\medskip

\item
При подключении параллельных линий к общим шинам на обоих концах НВЧЗ может сработать неселективно из-за реверса мощности.
В начальный момент повреждения на параллельной линии (линия W2) срабатывают блокирующие ИО полукомплектов А1 и А2, обеспечивая ускоренный запуск ВЧ-передатчиков.
Затем срабатывают отключающие ИО полукомплекта А2, но не срабатывают отключающие ИО полукомплекта А1, его ВЧ-передатчик остается в работе – наличие ВЧ-сигнала в канале связи предотвращает срабатывание НВЧЗ (см. рисунок \ref{pict_nvchz:REV_S_a}). 

\medskip

\begin{figure}[H]
\centering
\begin{subfigure}{0.6\textwidth}
    \centering
    \includegraphics[width=\textwidth]{\fbpath/561.0771 НВЧЗ/img/REV_S_a.pdf}
    \caption{\ начальный момент повреждения}
    \label{pict_nvchz:REV_S_a}
\end{subfigure}
\hfill
\begin{subfigure}{0.6\textwidth}
    \centering
    \includegraphics[width=\textwidth]{\fbpath/561.0771 НВЧЗ/img/REV_S_b.pdf}
	\caption{\ каскадное отключение параллельной линии}
    \label{pict_nvchz:REV_S_b}
\end{subfigure}
\caption{Работа НВЧЗ в режиме реверса мощности}
\end{figure}

\medskip

\item
В устройстве также предусматривается возможность ручного пуска ВЧ-передатчика через входной сигнал ФСУ \textit{«Ручной пуск ВЧ»}.
Выдаваемый ВЧ-сигнал при этом является манипулированным.
Ручной пуск разрешается, если в течение последних 600 мс не было срабатываний блокирующих органов или отключающего реле сопротивления.

\item
Кроме того, предусматривается возможность выдачи сплошного ВЧ-сигнала (логический сигнал \textit{«Пуск ВЧ от ДФЗ»}) в следующих случаях:
\begin{itemize} 
\item при срабатывании блокировки при сквозных токах через ошиновку, если введен программный ключ \textbf{«БСТО ДФЗ»};
\item от внешней команды пуска ВЧ-передатчика (входной сигнал ФСУ \textit{«Внешний пуск ВЧ»});
\item при обнаружении неисправности канала связи устройством автоматической проверки канала (АПК), если введен программный ключ \textbf{«Блок. ДФЗ от АПК»};
\item при обнаружении неисправности приемопередатчика, если введен программный ключ \textbf{«Блок. ДФЗ от Неиспр. ПП»};
\item при обнаружении неисправности терминала системой самодиагностики, если введен программный ключ \textbf{«Пуск ВЧ при неиспр. терм.»};
\item при выводе терминала или функции ДФЗ из работы, если введен программный ключ \textbf{«Пуск ВЧ при выводе»}.
\end{itemize}

\item
Выдача сплошного ВЧ-сигнала при срабатывании БСТО и от внешней команды пуска дополнительно блокируется при наличии сигнала останова ВЧ-передатчика (логический сигнал \textit{«Запрет пуска ВЧ от ДФЗ»}).
А при неисправности терминала, неисправности канала связи или приемопередатчика, а также при выводе терминала или функции ДФЗ из работы пуск ВЧ-передатчика никак не блокируется.

\item
Устройство АПК, входящее в состав ВЧПП, используется для автоматической проверки ВЧ-канала связи.
К устройству от ВЧПП подводятся сигналы \textit{«Контакт АПК»} и \textit{«Неиспр. ВЧПП»}.
Наличие сигнала \textit{«Контакт АПК»}, свидетельствует об отсутствии неисправности в канале связи.
При выявлении неисправности канала связи формируется логический сигнал \textit{«Сигн. неиспр. КС»}.
При ручном пуске ВЧ-передатчика, срабатывании блокирующих органов или отключающего реле сопротивления формируется сигнал \textit{«Запрет АПК от ДФЗ»}, блокирующий действие АПК.
Действие АПК также может быть выведено оперативно от внешнего сигнала \textit{«Вывод АПК»}.
При обнаружении неисправности приемопередатчика формируется логический сигнал \textit{«Сигн. Неиспр. ПП»}.
При введенных положениях программных ключей \textbf{«Блок. ДФЗ от АПК»} и \textbf{«Блок. ДФЗ от Неиспр. ПП»} дополнительно формируется сигнал \textit{«Блок. ДФЗ от ВЧПП»}, который действует на вывод ДФЗ из работы и пуск ВЧ-передатчика.

\item
Для сигнализации о наличии ВЧ-сигнала в канале связи реализован орган \textit{«ОСФ-Вызов»}, который вводится в работу при отсутствии сигнала \textit{«Запрет АПК»}.
Время срабатывания и возврата \textit{«ОСФ-вызов»} не превышает 20 мс.
Обеспечивается формирование сигнала \textit{«Вызов от ДФЗ»}, если сигнал срабатывания \textit{«ОСФ-вызов»} присутствует непрерывно в течение 5 с.

\medskip

\begin{figure}[H]
\centering
\includegraphics[width=1\textwidth,height=1\textheight,keepaspectratio]{\fbpath/561.0761 ДФЗ/img/Ctrl VCh DFZ.pdf}
\captionof{figure}{Функциональная схема контроля ВЧ-канала связи}\label{pict_dfz:Ctrl VCh DFZ}
\end{figure}

\medskip

\end{enumerate}


%%%%%%%%%%%%%%%%%%%%%%%%%%%%%%%%%%%%%%%%%%%%%%%%%% ОТКЛЮЧАЮЩИЕ ИО %%%%%%%%%%%%%%%%%%%%%%%%%%%%%%%%%%%%%%%%%%%%%%%%%%
\item Отключающие измерительные органы

\begin{enumerate}[label=\arabic{section}.\arabic{subsection}.\arabic{enumi}.\arabic*, labelsep=4pt, leftmargin=0em, itemindent=65pt, parsep=0pt]

\item
Подготовка цепи отключения осуществляется от отключающих органов.
Отключающие органы реализованы с возможностью независимой настройки их уставок по:
\begin{itemize} 
\item линейному току (модулю разности фазных токов) \textit{«ИО Iл откл ДФЗ»};
\item приращению тока прямой последовательности \textit{«ИО dI1 откл ДФЗ»};
\item приращению тока обратной последовательности \textit{«ИО dI2 откл ДФЗ»};
\item току обратной последовательности \textit{«ИО I2 откл ДФЗ»};
\item току нулевой последовательности \textit{«ИО 3I0 откл ДФЗ»};
\item по сопротивлению \textit{«РС откл ДФЗ»}.
\end{itemize}

\item
Отключающий орган по сопротивлению (отключающее РС) необходим для подхвата кратковременного срабатывания отключающих органов и обеспечения корректной работы ДФЗ при симметричных КЗ на защищаемой линии. 
Подхват вводится на время, равное 200 мс. Для сохранения селективности защиты при отключении внешнего КЗ (при неодновременной остановке работы ВЧПП) предусмотрено ограничение длительности срабатывания отключающего РС до 200 мс. 

\item
Отключающее РС имеет направленную характеристику срабатывания с охватом начала координат, которая задается уставками \textbf{«Rср РС откл ДФЗ»}, \textbf{«Хср РС откл ДФЗ»} и \textbf{«Ф1 РС откл ДФЗ»}. 
Для отстройки от нагрузки предусмотрен вырез в характеристике срабатывания, который определяется уставками \textbf{«Rнг РС откл ДФЗ»} и \textbf{«Фнг РС откл ДФЗ»}. 
Работа отключающего РС может быть заблокирована при выявлении неисправности в цепях напряжения при введенном программном ключе \textbf{«БНН РС откл ДФЗ»}.

\medskip

\begin{figure}[H]
\centering
\includegraphics[width=0.5\textwidth,height=0.5\textheight,keepaspectratio]{\fbpath/561.0761 ДФЗ/img/RS_DFZ.pdf}
\captionof{figure}{ Характеристика срабатывания отключающего РС ДФЗ}\label{pict_dfz:RS_DFZ}
\end{figure}

\medskip

\item
Для предотвращения ложного срабатывания ДФЗ при КЗ за трансформатором ответвления предусмотрен функциональный блок отстройки от повреждений за трансформаторами ответвлений (ОПТО). 
При введенном программном ключе \textbf{«Контр. ДФЗ от ОПТО»} подготовка цепи отключения разрешается только при наличии подтверждения о КЗ на защищаемой линии от ОПТО.

\item
При срабатывании отключающих органов и наличии подтверждения о КЗ на защищаемой линии от ОПТО формируется логический сигнал \textit{«Разрешение ОСФ»}, отвечающий за пуск ОСФ. 
Подключение ОСФ осуществляется с выдержкой времени: 10 мс для линий без ответвлений и 20 мс — при их наличии. 

\item
ОСФ предназначен для определения места повреждения по сдвигу фаз между ВЧ-пакетами, посылаемыми ВЧ-передатчиками обоих концов линии, то есть в конечном счёте по сдвигу фаз токов манипуляции.
Сравнение фаз осуществляется косвенно – по длительности пауз в ВЧ-сигнале.
При внешнем КЗ токи манипуляции находятся в противофазе, ВЧ-передатчики работают в чередующиеся полупериоды промышленной частоты, формируя непрерывный (сплошной) сигнал в ВЧ-канале, при котором ОСФ не срабатывает.
При внутреннем КЗ токи манипуляции на обоих концах совпадают по фазе, ВЧ-передатчики работают синхронно, что приводит к возникновению пауз в ВЧ-сигнале и срабатыванию ОСФ. 

\medskip

\begin{figure}[H]
\centering
\includegraphics[width=0.8\textwidth,height=0.8\textheight,keepaspectratio]{\fbpath/561.0761 ДФЗ/img/Action_DFZ.pdf}
\captionof{figure}{Принцип действия ДФЗ}
\label{pict_dfz:Action_DFZ}
\end{figure}

\medskip

\item
Длительность пауз, в принимаемом ВЧ-сигнале прямо зависит от величины угла сдвига фаз между токами манипуляции: чем меньше угол, тем длиннее паузы.
Уставками ОСФ задается величина сдвига фаз между токами манипуляции, при которой защита блокируется или срабатывает.
Предельное отклонение угла сдвига фаз от 180°, при котором защита ещё блокируется, называется углом блокировки.

\item
Срабатывание ОСФ происходит при превышении разности фаз уставки угла блокировки.
ОСФ выполнен трёхступенчатым, что позволяет реализовать зависимость времени срабатывания защиты от величины сдвига фаз.
Первая ступень ОСФ срабатывает при фиксировании одной паузы в ВЧ-сигнале, превышающей уставку угла блокировки первой ступени.
Вторая ступень – при двух последовательных превышениях уставки второй ступени.
Третья ступень – при трех последовательных превышениях уставки третьей ступени.
Углы блокировки задаются уставками \textbf{«Угол блок. ОСФ-1»}, \textbf{«Угол блок. ОСФ-2»} и \textbf{«Угол блок. ОСФ-3»} соответственно.

\item
Для компенсации расширения зоны блокировки, вызванного особенностями формирования сигнала \textit{«ВЧ ПРМ»} в ВЧПП используется корректировка уставок углов блокировки ОСФ на величину компенсации, задаваемую уставкой \textbf{«Удлинение ВЧ»}.

\medskip

\begin{figure}[H]
\centering
\includegraphics[width=0.6\textwidth,height=0.6\textheight,keepaspectratio]{\fbpath/561.0761 ДФЗ/img/OSF.pdf}
\captionof{figure}{Характеристика срабатывания ОСФ}
\label{pict_dfz:OSF}
\end{figure}

\medskip

\item
После срабатывания ОСФ осуществляется действие ДФЗ на отключение выключателя с регулируемой выдержкой времени \textbf{«Тср ДФЗ»}.
При одновременном срабатывании отключающих органов и появлении сигнала останова ВЧ-передатчика (логический сигнал \textit{«Запрет пуска ВЧ от ДФЗ»}) ОСФ шунтируется, и защита действует на отключение выключателя без выдержки времени.

\item
При подключении параллельных линий к общим шинам на обоих концах защита может сработать неселективно из-за реверса мощности.
При повреждении на параллельной линии срабатывают блокирующие органы полукомплектов, приводящие к выдаче ВЧ-передатчиком манипулированного ВЧ-сигнала.
Одновременно срабатывают и отключающие органы, запускающие ОСФ.
Однако срабатывания ОСФ и отключения неповрежденной линии не происходит, так как в данном режиме в ВЧ-канале присутствует сплошной сигнал.
При каскадном отключении одного из выключателей параллельной линии возникает реверс мощности.
Во время переходного процесса могут появиться паузы в ВЧ-сигнале, длительность которых достаточна для срабатывания ОСФ, что приведет к отключению неповрежденной линии.
Для исключения такого режима реализована блокировка при реверсе мощности: при одновременном наличии сигнала от отключающих органов и отсутствии срабатывания ОСФ в течение времени, задаваемого уставкой \textbf{«Тср блок. ДФЗ»}, срабатывает блокировка, которая продлевается на время, регулируемое уставкой \textbf{«Тпрод блок. ДФЗ»}.    

\item
Действие ДФЗ на отключение также блокируется при следующих условиях:
\begin{itemize} 
\item при срабатывании блокировки при сквозных токах через ошиновку, которая вводится программным ключом \textbf{«БСТО ДФЗ»};
\item при выводе терминала или функции ДФЗ из работы;
\item при обнаружении неисправности канала связи устройством автоматической проверки канала (АПК), если введен программный ключ \textbf{«Блок. НВЧЗ от АПК»};
\item при обнаружении неисправности приемопередатчика, если введен программный ключ \textbf{«Блок. НВЧЗ от Неиспр. ПП»};
\item при переводе действия ДФЗ только на сигнал (от внешнего сигнала \textit{«Перевод ДФЗ на сигнал»}).
\end{itemize}

\item
В сетях с тяговой нагрузкой, где наблюдаются значительные небалансы токов обратной последовательности, рекомендуется применять блокирующие и отключающие органы по приращениям токов прямой и обратной последовательностей.

\item
Время срабатывания защиты (без учета срабатывания выходного реле и с учетом задержки на подключение ОСФ) на отключение на двухконцевых линиях при кратности соответствующих величин, равной 3, не более 35 мс.

\medskip

\begin{figure}[H]
\centering
\includegraphics[width=1\textwidth,height=1\textheight,keepaspectratio]{\fbpath/561.0761 ДФЗ/img/DFZ.pdf}
\captionof{figure}{Функциональная схема ДФЗ}
\label{pict_dfz:DFZ}
\end{figure}

\medskip

\begin{figure}[H]
\centering
\includegraphics[width=1\textwidth,height=1\textheight,keepaspectratio]{\fbpath/561.0761 ДФЗ/img/VChPP.pdf}
\captionof{figure}{Функциональная схема управления ВЧПП (ДФЗ)}
\label{pict_dfz:VChPP}
\end{figure}

\medskip

\begin{figure}[H]
\centering
\includegraphics[width=1\textwidth,height=1\textheight,keepaspectratio]{\fbpath/561.0761 ДФЗ/img/IO_DFZ.pdf}
\captionof{figure}{Измерительные органы ДФЗ}
\label{pict_dfz:IO_DFZ}
\end{figure}

\medskip

\item
Уставки ДФЗ представлены в таблице \ref{desc_dfz:tbl1}.

\small
\begin{longtable}{|>{\centering\arraybackslash}m{5.3cm}|>{\centering\arraybackslash}m{3.3cm}|>{\centering\arraybackslash}m{4.2cm}|>{\centering\arraybackslash}m{1.8cm}|>{\centering\arraybackslash}m{1cm}|}
% Первая страница: полный заголовок
\caption{Уставки ДФЗ\hfill\vspace{-0.5\baselineskip}}\label{desc_dz:tbl1}\\
\hline
\rowcolor{gray!30}
Параметр & Условное обозначение на схеме & Значение/ Диапазон & Единица измерения & Шаг \\
\hline
\endfirsthead

% Промежуточные страницы: «Продолжение таблицы»
\caption*{\hspace{3pt}\emph{Продолжение таблицы \ref{desc_dz:tbl1}}\hfill\ }\vspace{-0.5\baselineskip} \\
\hline
\rowcolor{gray!30}
Параметр & Условное обозначение на схеме & Значение/ Диапазон & Единица измерения & Шаг \\
\hline
\endhead

%===m>ДФЗ|-
%===m
\end{longtable}
\normalsize

\end{enumerate}

\end{enumerate}
